\documentclass[a4paper,11pt,fleqn]{scrartcl}%fleqn pour aligner les equations à gauche
\usepackage{../../Styles/Eval}
\usepackage{longtable}

\newcommand{\chnum}{Chapitre 5 et 18}
\newcommand{\chtitle}{Mouvement des planètes et lunette astronomique}
\newcommand{\dtype}{DS 5}

\begin{document}
\begin{center}
\renewcommand{\arraystretch}{1.6}

\section*{Correction}

\renewcommand{\arraystretch}{1.6}

\begin{longtable}{||p{.1\linewidth}|p{.78\linewidth}|p{.07\linewidth}|}
\hline
\textbf{Question} & \textbf{Réponse} & \textbf{Barème} \\
\hline
\endfirsthead

\hline
\textbf{Question} & \textbf{Réponse} & \textbf{Barème} \\
\hline
\endhead

\hline
\multicolumn{3}{r}{Suite page suivante} \\
\hline
\endfoot

\hline
\textbf{Total} & & \textbf{20 pts} \\
\hline
\endlastfoot

1 &
\textbf{Première loi de Kepler :} Les planètes décrivent des orbites elliptiques dont le Soleil occupe l'un des foyers. 

\textbf{Deuxième loi de Kepler :} Le segment reliant la planète au Soleil balaie des aires égales pendant des durées égales. 
& 2 pts \\

\hline

2 &
La deuxième loi indique que des aires égales sont balayées pendant des durées égales.  
Lorsque Mars est plus proche du Soleil (périhélie), elle parcourt une portion d'orbite plus grande pendant la même durée que lorsqu'elle est plus éloignée (aphélie).  

La vitesse de Mars varie donc au cours du temps : le mouvement n'est pas uniforme.
& 2 pts \\

\hline

3 &
Ligne 8 du programme :

\begin{verbatim}
T2_a3.append(T[i]**2/a[i]**3)
\end{verbatim}

Cette ligne calcule le rapport $T^2/a^3$ pour chaque planète.
& 2 pts \\

\hline

4 &
Les lignes 16 et 17 permettent de calculer puis d'afficher la valeur moyenne du rapport $T^2/a^3$ pour les planètes étudiées.

La finalité est de vérifier que ce rapport est constant pour toutes les planètes.

Unités :$$
\frac{T^2}{a^3} \quad \Rightarrow \quad \text{s}^2\cdot \text{m}^{-3}$$

& 2 pts \\

\hline

5 &
La représentation graphique montre que les points expérimentaux sont alignés sur une droite.

Cela signifie que le rapport $T^2/a^3$ est constant pour les planètes du système solaire.

Cela confirme la \textbf{troisième loi de Kepler} :$$\frac{T^2}{a^3}= \text{constante}$$

& 2 pts \\

\hline

6 &
Période de Mars : $
T_{Mars}=687\,\text{jours}$

Conversion en secondes : 
$
T_{Mars}=687\times 86400
$

Troisième loi de Kepler : $\frac{T^2}{a^3}= \text{constante}$

En utilisant les données de la Terre : 
\[
\frac{T_{Mars}^2}{a_{Mars}^3}=\frac{T_T^2}{a_T^3}
\]

On obtient : $a_{Mars} \approx 2,28\times 10^{11}\,m$

soit environ :
\[
a_{Mars} \approx 1,52 \, UA
\]

Cette distance est supérieure à celle de la Terre (1 UA) mais inférieure à celle de Jupiter, ce qui correspond à la \textbf{4\textsuperscript{ième} planète à partir du Soleil}.

& 4 pts \\

\hline

7 &
Dans une lunette astronomique :

\begin{itemize*}
\item $L_1$ est l'\textbf{objectif} (lentille proche de l'objet observé)
\item $L_2$ est l'\textbf{oculaire} (lentille proche de l'œil)
\end{itemize*}

& 1 pt \\

\hline

8 &
Les rayons lumineux issus de $B_\infty$ arrivent parallèles sur l'objectif $L_1$.  

Ils convergent vers le foyer image $F'_1$ de l'objectif et forment l'image réelle intermédiaire $A_1B_1$.  

Cette image se trouve dans le plan focal objet de l'oculaire $L_2$.
L'oculaire transforme alors les rayons en un faisceau parallèle permettant l'observation à l'infini.
& 2 pts \\

\hline

9 &
$\theta_1$ : angle sous lequel Mars est vue à l'œil nu.  

$\theta_2$ : angle sous lequel Mars est vue à travers la lunette.  

Ces deux angles sont mesurés depuis l'œil de l'observateur, l'angle $\theta_2$ étant plus grand que $\theta_1$.

& 1 pt \\

\hline

10 &
Le grossissement est défini par :
$G=\frac{\theta_2}{\theta_1}$

Dans une lunette astronomique afocale : $G=\frac{f'_{obj}}{f'_{ocu}}$

où :
\begin{itemize*}
\item $f'_{obj}$ est la distance focale de l'objectif
\item $f'_{ocu}$ est la distance focale de l'oculaire
\end{itemize*}

& 2 pts \\

\hline

11 &
Angle sous lequel Mars est vue depuis la Terre : $
\theta_1 \approx \frac{h}{d}$

avec : $
h=6794\,km$
\[
d=62,07\times10^6\,km
\]
\[
\theta_1 \approx 1,1\times10^{-4}\,rad
\]

Pour voir Mars aussi grosse que la Lune :
\[
\theta_2 \ge 9,0\times10^{-3}\,rad
\]

Grossissement minimal :
$$G_{min}=\frac{\theta_2}{\theta_1} \rightarrow G_{min}\approx 82$$
Calcul du grossissement de la lunette :
\[
G=\frac{f'_{obj}}{f'_{ocu}}
\]
avec $f'_{obj}=900\,mm$
\[
G_{10}=\frac{900}{10}=90
\]
\[
G_{25}=36
\]
\[
G_{40}=22,5
\]
Seul l’oculaire de \textbf{10 mm} permet d’obtenir un grossissement suffisant pour voir Mars au moins aussi grosse que la Lune vue à l’œil nu.
& 4 pts \\
\hline

\end{longtable}
\end{center}

\end{document}